\subsection{Does the correction matter on sparser logs?}\label{sec:sparse}

\begin{table}[t]\centering\small
\caption{KuaiRec with the log thinned to a fraction $q$ of its exposures
(nuisances re-fitted, additive imputation, $\tau\in\{0.05,0.01\}$, $\lambda\in\{0,1\}$,
selected on validation). nDCG@20 of the uncorrected and corrected operators, and
the lowest top-20 overlap and Spearman correlation between the two rankings at
matched configurations.}\label{tab:sparse}
\begin{adjustbox}{max width=\textwidth}
\begin{tabular}{cccccccc}
\toprule
thinning $q$ & density & DR & DRUP & DR 5 hops & DRUP 5 hops & top-20 overlap & Spearman\\
\midrule
1 & 13.38\% & 0.6302 & 0.6302 & 0.6306 & 0.6306 & 0.971 & 0.999\\
0.1 & 1.34\% & 0.6097 & 0.6102 & 0.6109 & 0.6111 & 0.935 & 0.999\\
0.03 & 0.40\% & 0.6073 & 0.6069 & 0.6085 & 0.6081 & 0.985 & 1.000\\
\bottomrule
\end{tabular}

\end{adjustbox}
\end{table}

Theorem~\ref{thm:lower} predicts a larger bias when propensities are small. We
therefore thinned the KuaiRec log to 10\% and 3\% of its exposures, which lowers the
median logged propensity from 0.48 to 0.045 and 0.013, re-fitted the nuisances
and re-selected every operator (Table~\ref{tab:sparse}). The corrected and the
uncorrected operators still rank almost identically: their top-20 lists overlap
by at least 93\%, the rank correlation is at least 0.999, and the accuracy
differences stay below 0.0005 nDCG@20 in either direction, at three and at five
hops. On these data the repeated-walk term is small compared with the score
differences that decide the ranking, and the value of the correction lies in
the properties of the scores, not in the rankings they induce.
